\section{Especificación vigente de indicadores}
\label{anexo:a}

La especificación reproducible se mantiene en \texttt{resultados/base\_trabajo/especificacion\_indicadores.csv}. La auditoría estadística registra 17 definiciones operativas: siete variables de una matriz candidata provisional y otras variables descriptivas, contextuales o pendientes de revisión. La tabla resume únicamente las siete candidatas; su presencia no implica aprobación conceptual ni que se haya ejecutado el PCA.

\begin{longtable}{@{}p{0.31\textwidth}p{0.27\textwidth}p{0.32\textwidth}@{}}
\caption{Variables de la matriz candidata y principales reservas conceptuales}\label{tab:especificacion_indicadores}\\
\toprule
\textbf{Variable} & \textbf{Medida observada} & \textbf{Reserva para interpretación} \\
\midrule
\endfirsthead
\toprule
\textbf{Variable} & \textbf{Medida observada} & \textbf{Reserva para interpretación} \\
\midrule
\endhead
Actividad combinada & Proporción de población 10+ en categorías combinadas de actividad & Agrupa búsqueda de empleo, oficios del hogar, incapacidad permanente y otras categorías; no equivale a desempleo ni a una medida unidimensional de vulnerabilidad. \\
Sin energía eléctrica & Viviendas ocupadas sin energía eléctrica & Indicador de servicio básico; interpretar con contexto de oferta territorial. \\
Sin gas de red & Viviendas ocupadas sin gas de red pública & La disponibilidad depende de infraestructura y oferta local. \\
Sin recolección de basuras & Viviendas ocupadas sin recolección de basuras & La cobertura puede reflejar provisión del servicio y dispersión poblacional. \\
Sin internet & Viviendas ocupadas sin internet & Distribución próxima al techo en parte del universo; revisar capacidad discriminante y no respuesta. \\
Hogares con menores de 15 & Proporción de hogares con menores de 15 años & Es una característica de composición del hogar, no una privación directa. \\
Estratos 1 y 2 & Proporción de personas en estratos 1 y 2 & Proxy administrativo con cobertura territorial desigual; MSA individual reportado de 0,414. \\
\bottomrule
\end{longtable}

Los diagnósticos disponibles son preliminares (KMO global 0,642; Bartlett $\chi^2=4664,873$, 21 grados de libertad, $p<0,001$). No certifican por sí mismos la pertinencia del análisis ni sustituyen la revisión conceptual, la evaluación de estabilidad y la decisión sobre balance entre dominios. El IPM Municipal se reserva para análisis descriptivos y eventual contraste externo; cualquier asociación debe interpretarse considerando el solapamiento conceptual de sus privaciones con variables censales.

